% problem-000246 7 von Richter重排机理
\section*{7. von Richter重排机理}
1871年发现的vonRichter重排反应（如下图）是有机机理研究领域中非常有趣的一个反应，其反应机理被反复推倒、重建，历经百年才被逐步证实。

（图：）

早期化学家提出了一系列简单的反应机理，但相继被否定。上世纪50年代，通过一系列实验研究，Bunnett 和Rauhut 提出了如下机理

（图：）
第11页 共17页

第38届中国化学奥林匹克（决赛）理论第一场试题 2024年10月26日广州

\subsection*{7-11960}
年的实验研究表明，反应的含N产物主要为N2，从而对上述机理提出了质 $\Nu _ { 2 } ,$ 疑。尽管有人辩解称，上述机理中A水解产生的NH3可被后面解离的 $\mathrm { N O } _ { 2 } ^ { - }$ 或活性中间体 B 氧化生成 $\aleph _ { 2 } ,$ 但该机理还是被当年的 $^ { 1 5 } \mathsf { N }$ 同位素示踪实验所否定。用何种15N同位素试剂可实现上述目的？

\subsection*{7-2}
目前确立下来 von Richter 重排的部分机理如下

（图：）

其中D和E为稳定的中间体，可由其他途径制备后在相同条件下继续反应得到vonRichter产物，从而有力地证实了该机理。已知D为亚硝基化合物，E可被环戊二烯捕获得到相对分子量为198.1的化合物F。写出D、E和F的结构式。

\subsection*{7-3}
硝基芳香族化合物BTZ043是潜在的结核病治疗药物。研究表明，BTZ043在亲核试剂作用下产生含亚硝基的中间体可能是其生物活性的关键。

为模拟 BTZ043与酶中半胱氨酸残基的作用，在 $\tt p H 1 1 \sim 1 2$ 的乙腈-水溶剂中用Boc-cysteamine（N-Boc 氨基乙硫醇）对 BTZ043 进行处理，并加入1.3-环己二烯捕获反应产生的中间产物，液相色谱-质谱（LC-MS）可检测到两组[M+H]的分子离子峰，m/z值分别为：496.2与353.2。画出G和H的结构式。

（图：）

\subsection*{7-4}
为进一步探讨亲核试剂对 BTZ043的作用，在乙腈-水混合溶剂中用 KCN 处理BTZ043，分离得到两种主要产物J和K，并在过程中产生另一种相对分子量为456.1的物质L。研究表明，生成J和K要经历共同的关键中间体M $( \mathbf { C } _ { 1 8 } \mathbf { H } _ { 1 7 } \mathbf { F } _ { 3 } \mathbf { N } _ { 4 } 0 _ { 5 } \mathbf { S } )$ 。

（图：）
第12页 共17页

第38届中国化学奥林匹克（决赛）理论第一场试题2024年10月26日广州

\subsection*{7-4}
-1画出L、M的结构式。

\subsection*{7-4}
-2画出产生L过程中的关键中间体。

\subsection*{7-4}
-3 画出由M转化为J过程中的关键中间体。

\subsection*{7-4}
-4 画出由M 转化为 K 过程中的关键中间体。
